\begin{lemma}
\label{editorial-source-lemma-image-finite-type-map-Jacobson-rings}
Let $R \to S$ be a finite type ring map of Jacobson rings.
Denote $X = \Spec(R)$ and $Y = \Spec(S)$.
Write $f : Y \to X$ for the induced
map of spectra. Let $E \subset Y = \Spec(S)$ be a
constructible set. Denote with a subscript ${}_0$ the set
of closed points of a topological space.
\begin{enumerate}
\item We have $f(E)_0 = f(E_0) = X_0 \cap f(E)$.
\item A point $\xi \in X$ is in $f(E)$ if and only if
$\overline{\{\xi\}} \cap f(E_0)$ is dense in $\overline{\{\xi\}}$.
\end{enumerate}
\end{lemma}

\begin{proof}
We have a commutative diagram of continuous maps
$$
\xymatrix{
E \ar[r] \ar[d] & Y \ar[d] \\
f(E) \ar[r] & X
}
$$
Suppose $x \in f(E)$ is closed in $f(E)$. Then $f^{-1}(\{x\})\cap E$
is nonempty and closed in $E$. Applying
Topology, Lemma \ref{topology-lemma-jacobson-inherited}
to both inclusions
$$
f^{-1}(\{x\}) \cap E \subset E \subset Y
$$
we find there exists a point $y \in f^{-1}(\{x\}) \cap E$ which is
closed in $Y$. In other words, there exists $y \in Y_0$ and $y \in E_0$
mapping to $x$. Hence $x \in f(E_0)$.
This proves that $f(E)_0 \subset f(E_0)$.
Proposition \ref{editorial-source-proposition-Jacobson-permanence} implies that
$f(E_0) \subset X_0 \cap f(E)$. The inclusion
$X_0 \cap f(E) \subset f(E)_0$ is trivial. This proves the
first assertion.

\medskip\noindent
Suppose that $\xi \in f(E)$. According to
Lemma \ref{editorial-source-lemma-characterize-image-finite-type}
the set $f(E) \cap \overline{\{\xi\}}$ contains a dense
open subset of $\overline{\{\xi\}}$. Since $X$ is Jacobson
we conclude that $f(E) \cap \overline{\{\xi\}}$ contains a
dense set of closed points, see Topology,
Lemma \ref{topology-lemma-jacobson-inherited}.
We conclude by part (1) of the lemma.

\medskip\noindent
Conversely, suppose $\overline{\{\xi\}}\cap f(E_0)$ is dense
in $\overline{\{\xi\}}$. By Lemma \ref{editorial-source-lemma-constructible-is-image}
choose a finitely presented map $S\to S'$ whose spectrum map
$g:Y'=\Spec(S')\to Y$ has image exactly $E$.
The ring $S'$ is Jacobson by Proposition \ref{editorial-source-proposition-Jacobson-permanence}.
Part (1) of the present lemma, already proved and applied to $S\to S'$,
gives $g(Y'_0)=E_0$. Put $h=f\circ g$, so
$h(Y'_0)=f(E_0)$ and $h(Y')=f(E)$, retaining the original $S$ and $E$.
If $\xi$ corresponds to $\mathfrak p\subset R$, use the exact diagram
$$
\xymatrix{
S'\ar[r] & S'/\mathfrak pS'\\
R\ar[r]\ar[u] & R/\mathfrak p\ar[u]
}
$$
Every point of $Y'_0$ whose image lies in $V(\mathfrak p)$ corresponds
to a prime of $S'$ containing $\mathfrak pS'$, hence to a point of
$\Spec(S'/\mathfrak pS')$. The spectrum image of
$R/\mathfrak p\to S'/\mathfrak pS'$ therefore contains the given
dense subset $h(Y'_0)\cap V(\mathfrak p)$ of
$V(\mathfrak p)=\Spec(R/\mathfrak p)$.
By Lemma \ref{algebra-lemma-domain-image-dense-set-points-generic-point},
there is a prime $\overline{\mathfrak q}'$ of $S'/\mathfrak pS'$
contracting to $(0)$ in $R/\mathfrak p$. Its inverse image
$\mathfrak q'\subset S'$ contracts to $\mathfrak p$ in $R$.
The point $g(\mathfrak q')$ belongs to $E$ and its image under the
original $f$ is $\xi$. Consequently $\xi\in f(E)$, as required.
\end{proof}